\documentclass[a4paper]{article}

% --- ESSENTIAL PACKAGES FOR COMPACTION ---
\usepackage[a4paper, margin=0.2in]{geometry} % set tight margins
\usepackage{multicol}      % create a multi-column layout
\usepackage{titlesec}      % customize section/subsection titles
\usepackage{paralist}      % creating compact lists
\usepackage{amsmath}       % math environments

% --- UTILITY & STYLING PACKAGES ---
\usepackage{xcolor}        % custom colors
\usepackage{graphicx}      % images
\usepackage{listings}      % code snippets
\usepackage{hyperref}      % clickable links

% --- COLOR DEFINITIONS ---
\definecolor{codegray}{rgb}{0.9,0.9,0.9} % code background
\definecolor{darkblue}{rgb}{0.0,0.0,0.5} % for comments
\definecolor{defbox_bg}{rgb}{0.9, 0.9, 1.0} % definition box background
\definecolor{defbox_fr}{rgb}{0.3, 0.3, 0.7} % frame

% --- GLOBAL LAYOUT & SPACING CONFIGURATION ---

% set paragraph spacing: no indent, with tiny vertical gap between paragraphs
\setlength{\parindent}{0pt}
\setlength{\parskip}{0.5ex}

% make section and subsection titles extremely compact
\titlespacing*{\section}{0pt}{1.2ex}{0.8ex}
\titlespacing*{\subsection}{0pt}{1ex}{0.5ex}
\titleformat*{\section}{\bfseries\large}
\titleformat*{\subsection}{\bfseries\normalsize}

% make lists as compact as possible
\let\itemize\compactitem
\let\enditemize\endcompactitem
\let\enumerate\compactenum
\let\endenumerate\endcompactenum

% --- CODE LISTING STYLE ---
\lstdefinestyle{mystyle}{
    backgroundcolor=\color{codegray},   
    commentstyle=\color{darkblue},
    keywordstyle=\color{blue},
    numberstyle=\tiny\color{gray},
    stringstyle=\color{red},
    basicstyle=\ttfamily\tiny,
    breaklines=true,
    numbers=left,                    
    numbersep=4pt,                  
    showspaces=false,                
    showstringspaces=false,
    showtabs=false,                  
    tabsize=2
}
\lstset{style=mystyle}

% --- HYPERREF SETUP---
\hypersetup{
    colorlinks=true, linkcolor=blue, urlcolor=cyan,
}

% --- CUSTOM COMMANDS/ENVIRONMENTS ---
% tight box for definitions
\newenvironment{definition}[1]
    {\par\noindent\begin{tcolorbox}[colback=blue!5!white,colframe=blue!75!black,title=#1,fonttitle=\bfseries,boxsep=1mm,left=1mm,right=1mm,top=1mm,bottom=1mm]}
    {\end{tcolorbox}}

% simpler alternative w/o tcolorbox
% \newcommand{\definition}[2]{\par\noindent\fcolorbox{defbox_fr}{defbox_bg}{\begin{minipage}{0.95\columnwidth}\textbf{#1:} #2\end{minipage}}\par}

\begin{document}

% --- CUSTOM, MINIMAL TITLE ---
\noindent
\textbf{\tiny Cheat Sheet - Algoritmi Avansati}
\hrule
\vspace{0.25ex}

% --- GLOBAL SETTINGS FOR THE DOCUMENT BODY ---
\pagestyle{empty}
\footnotesize


\section{}

\begin{multicols}{2}

% \begin{definition}{titlu}
P (Polynomial) - Clasa de probleme pentru care le putem afla solutia in timp polinomial \\
Np (Nondeterministic Polynomial) - Clasa de probleme pentru care putem verifica in timp polinomial daca un rezultat este solutie corecta pentru problema noastra. \\
Problema de Optim: \\
Informal spus este problema in care trebuie sa gasesti o “cea mai buna” solutie/constructie fezabila. \\
“Cea mai buna” - poate avea doua sensuri:   \\
\begin{itemize}
\item Fie avem o problema de minimizare precum Problema Comis-voiajorului. \\
\item Fie o problema de maximizare precum cea de a găsi o acoperire de cardinal maxim pentru multimea varfurilor unui graf \\
\end{itemize}
Fie P - o problema de optim, și I o intrare pe aceasta problema. Vom nota cu OPT(I) ”valoarea” soluției optime. \\
În mod analog, atunci când propunem un algoritm care să ofere o soluție fezabilă pentru problema
noastră, vom nota ”valoarea” acelei soluții cu ALG(I). \\
De cele mai multe ori, atunci când nu se crează confuzie, vom simplifica notațiile folosind termenii
”OPT”, respectiv ”ALG” \\
Pe parcursul prezentării vom presupune că atât OPT, cât și ALG sunt $\geq0$. \\
Pentru a justifica ca un algoritm este util, aceste trebuie insotit de o justificare
ca solutia oferita este fezabila pentru problema, precum si o relatie intre ALG si OPT.
Acest tip de relatie este descrisa astfel: \\
Definitie 1: 
\begin{itemize}
\item un algoritm ALG pentru o problema de minimizare se numeste 
$\rho$-aproximativ, pentru o valoare $\rho >1$, daca ALG(I) $\leq \rho \cdot$OPT(I) pentru $\forall$ I - intrare \\
\item un algoritm ALG pentru o problema de maximizare se numeste 
$\rho$-aproximativ, pentru o valoare $\rho <1$, daca ALG(I) $\geq \rho \cdot$OPT(I) pentru $\forall$ I - intrare \\
\end{itemize}
Observatie pentru problemele de minim: Orice algoritm $\rho$-aproximativ este la randul lui
$\rho$'-aproximativ pentru orice $\rho$' > $\rho$. De aceea, in cazul unul algoritm ALG pentru o problema de minimizare, spre exemplu, trebuie 
ca justificarea ce insoteste pe ALG sa ofere cea mai mica valoare $\rho$ pentru care ALG este $\rho$-aproximativ.
Definitie 2: \\
Fie ALG un algoritm $\rho$-aproximativ pentru o problema de minimizare. Spunem ca 
factorul de aproximare este "tight bounded" atunci cand avem $\rho$ = $supremum_I$
$\frac{ALG(I)}{OPT(I)}$. Ca sa aratam ca un algoritm este $\rho$-aproximativ "tight bounded",
trebuie deci sa justificam urmatoarele 2 lucruri:
\begin{itemize}
    \item Trebuie sa aratam ca este $\rho$-aproximativ, adica ALG(I)$\leq \rho \cdot$OPT(I) pentru orice intrare I
    \item Pentru orice $\rho$'<$\rho$ exista un I pentru care ALG(I)>$\rho$'$\cdot$OPT(I).
     Adesea totusi ne este la indemana sa aratam ca exista un I pentru care ALG(I) = $\rho \cdot$OPT(I)

\end{itemize}

Problema rucsacului 1/0

Enunt pe scurt: Trebuie sa gasim o submutime de obiecte de valoare totala maxima, fara ca greutatea lor totala sa depaseasca o capacitate data a rucsacului. Obiectele sunt puse integral in rucsac sau sunt data deoparte. Nu pot fi fractionate. \\
In mod normal este programare dinamica cu complexiate n$\cdot$c. Presupunere: fiecare obiect are o greutate mai mica sau egala cu capacitatea rucsacului! \\

Rezolvare propusa: \\

Fie L - lista obiectlor sortate dupa raportul $\frac{valoare}{greutate}$ \\
Fie $ O_p$ - obiectul cu profitul cel mai mare din lista de obiecte. \\
S = 0, G = capacitatea rucsacului; \\
Pentru fiecare O:L daca greutatea(O)$leq$G, atunci S += val(O), G-= greutate(O) \\
ALG(I) = max(S, $O_p$) \\
Justificarea:
\begin{itemize}
    \item In primul rand, algoritmul de mai sus ne ofera o solutie fezabila. Elementele care au suma valorilor S vor avea o greutate totala $\geq$ capacitatea rucsacului, respectiv $O_p$ incape si el in rucsac de unul singur.
    \item Trebuie sa justificam doar factorul de aproximare. Fie $OPT_{1/0}$ valoarea optima pentru Problema Rucsacului in varianta 1/0, respectiv $OPT_G$ valoarea optima, furnizata de algoritmul de tip greedy pentru problema rucsacului in care avem voie sa fractionam obiecte pentru a le incarca in rucsac. 
    \item Cum este $OPT_{1/0}$ fata de $OPT_G$? \quad $OPT_{1/0} \geq OPT_G$
    \item Fie k indicele primului obiect care nu este adaugat in algoritmul de la inceputul paginii.
    \item $OPT_{1/0} \leq OPT_G \leq \sum_{1 \leq i \leq k} val(O_i) = \sum_{1 \leq i \leq k} val(O_i) + val(O_k) leq \sum_{1 \leq i \leq k} val(O_i) + val(O_p)$
    \item ALG = max(S, $O_P$)
    \item $OPT_{1/0} leq \sum_{1 \leq i \leq k} val(O_i) + val(O_p) leq ALG + ALG = 2 \cdot ALG. OPT_{1/0} \leq 2 \cdot ALG$
    \item Ex intrare pt care abaterea e maxima: G = 100; Ob (val/greutate) = [(50+eps1)/(50+eps2), 50/50, 50/50] cu eps1 > eps2 > 0
    \item Profitul maxim este 100, iar profitul solutiei algoritmului este 50+eps1. ALG(I)$ \cong \frac{1}{2} \cdot OPT(I) deci \frac{1}{2}$ este un "tight upper bound"
% \end{definition}
\end{itemize}
\end{multicols}

\section{}


\end{document}